%-------------------------------------------------------------------------------
%	SECTION TITLE
%-------------------------------------------------------------------------------
\cvsection{个人简介}
%\cvsection{Summary}


%-------------------------------------------------------------------------------
%	CONTENT
%-------------------------------------------------------------------------------
\begin{cvparagraph}

%---------------------------------------------------------
何洪津，男，生于1984年8月，籍贯湖南郴州，
宁波大学数学与统计学院教授、博士生导师，
担任中国运筹学会数学规划分会青年理事、
中国运筹学会算法软件与应用分会理事、
中国优选法统筹法与经济数学研究会-经济数学与管理数学分会常务理事。
2012年6月博士毕业于南京师范大学计算数学专业。
曾受邀访问澳大利亚科廷大学、香港理工大学、
香港浸会大学、台湾中山大学、台湾高雄医学大学。
主要研究兴趣为结构型优化、变分不等式和（张量特征值）互补问题及其在信号/图像处理、
统计学习、交通管理、机器学习等领域中的应用。
研究成果发表在Numerische Mathematik, Inverse Problems, 
Journal of Scientific Computing, Advances in Computational Mathematics, 
Journal of Optimization Theory and Applications, 
Computational Optimization and Applications 等国际权威期刊，
已发表学术论文80余篇（SCI期刊论文70余篇）。
已主持完成国家自然科学基金面上项目、青年基金项目、
浙江省自然科学基金一般探索（面上）项目共4 项，
现正主持国家自然科学基金面上项目1项、浙江省自然科学基金重点项目1项、
宁波市自然科学基金重点项目1项，参与完成浙江省自然科学基金重大、重点项目，
国家自然科学基金面上项目共5项。
2017年10月入选``浙江省高校中青年学科带头人"培养对象，
2017年12月入选杭州电子科技大学首届“优秀骨干教师支持计划”，
2025年入选浙江省高层次人才特支计划青年人才，
2026年入选宁波市领军与拔尖人才（第一层次）。
%万人计划青年拔尖人才培养对象。
担任美国数学会《数学评论》评论员，
多个优化和计算数学类学术期刊的审稿人。

\end{cvparagraph}
